\chapter{Laplace's Equation}

\section*{Introduction}

Laplace's equation, $\nabla^{2}\varphi = 0$, is one of the most ubiquitous partial differential equations in mathematical physics. Named after Pierre-Simon Laplace (though studied earlier by Daniel Bernoulli in 1738 and Hermann von Helmholtz in 1858), it describes the behavior of scalar potentials in regions free of sources: electrostatic potential in charge-free space, gravitational potential in vacuum, steady-state temperature distributions, and incompressible, irrotational fluid flow. Its solutions, called harmonic functions, possess remarkable properties: they are infinitely differentiable, satisfy the maximum principle, and can be constructed by separation of variables, conformal mapping, or the mean-value theorem.

Whenever a physical field is source-free and in equilibrium, the governing potential satisfies Laplace's equation. In electrostatics, $\nabla^{2}\varphi = 0$ in regions without charge; in gravitation, $\nabla^{2}\Phi = 4\pi G\rho$ reduces to $\nabla^{2}\Phi = 0$ in vacuum; in heat conduction, the steady-state temperature satisfies $\nabla^{2}T = 0$; and in fluid mechanics, the velocity potential of an irrotational, incompressible flow satisfies $\nabla^{2}\varphi = 0$. The universality of Laplace's equation arises from the linearity and isotropy of the underlying physics.


\begin{equation}
\boxed{\; \nabla^{2}\varphi = 0 \quad\Longleftrightarrow\quad \frac{\partial^{2}\varphi}{\partial x^{2}} + \frac{\partial^{2}\varphi}{\partial y^{2}} + \frac{\partial^{2}\varphi}{\partial z^{2}} = 0\;}
\label{eq:laplace}
\end{equation}

\section*{Fundamental Equations Used}

The derivation starts from Gauss's law in differential form, which states that the divergence of a conservative field vanishes in source-free regions.

\section*{Assumptions}

\textbf{Assumptions:} The domain is source-free ($\rho = 0$, or $q = 0$, or steady state); the medium is homogeneous and isotropic; boundary conditions are specified.


\textbf{Limitations:} Does not describe fields in the presence of sources (use Poisson's equation $\nabla^{2}\varphi = -\rho/\varepsilon_{0}$); does not describe time-dependent phenomena (use the heat or wave equation).


\section*{Mathematical Derivation}

\textbf{Step 1: Start from the divergence of a conservative field.}

Consider a region of space free of sources (no charge, no mass, no heat generation). If $\varphi$ is a scalar potential whose gradient gives a conservative field $\mathbf{E} = -\nabla\varphi$, then Gauss's law (or the divergence theorem applied to the flux) requires:
\begin{equation}
\nabla\cdot\mathbf{E} = 0 \quad\Longrightarrow\quad \nabla\cdot(-\nabla\varphi) = 0 \quad\Longrightarrow\quad \nabla^{2}\varphi = 0.
\end{equation}

\textbf{Step 2: Show the Laplacian is the sum of second partial derivatives.}

In Cartesian coordinates, $\nabla^{2}\varphi = \partial^{2}\varphi/\partial x^{2} + \partial^{2}\varphi/\partial y^{2} + \partial^{2}\varphi/\partial z^{2}$. The physical meaning is that the Laplacian at a point measures the difference between the value of $\varphi$ at that point and its average over a small sphere centered on the point.

\textbf{Step 3: Prove the mean-value property.}

Integrate $\nabla^{2}\varphi$ over a ball $B_R$ of radius $R$ centered at $\mathbf{r}_0$:
\begin{equation}
\int_{B_R} \nabla^{2}\varphi\,dV = \oint_{\partial B_R} \nabla\varphi\cdot\hat{n}\,dA = 0.
\end{equation}
By the divergence theorem, the volume integral of the Laplacian equals the flux of $\nabla\varphi$ through the sphere. Using the spherical symmetry of the surface integral:
\begin{equation}
\oint_{\partial B_R} \nabla\varphi\cdot\hat{n}\,dA = 4\pi R^{2}\,\langle\nabla\varphi\cdot\hat{n}\rangle_{\text{avg}} = 0.
\end{equation}
Integrating again with respect to $R$ yields the mean-value property:
\begin{equation}
\varphi(\mathbf{r}_0) = \frac{1}{4\pi R^{2}}\oint_{\partial B_R} \varphi(\mathbf{r})\,dA = \frac{1}{\frac{4}{3}\pi R^{3}}\int_{B_R} \varphi(\mathbf{r})\,dV.
\end{equation}

\textbf{Step 4: Prove the maximum principle.}

Suppose $\varphi$ has a local maximum at an interior point $\mathbf{r}_0$. Then by the mean-value property:
\begin{equation}
\varphi(\mathbf{r}_0) = \frac{1}{|\partial B_R|}\oint_{\partial B_R}\varphi\,dA.
\end{equation}
If $\varphi(\mathbf{r}_0)$ is a strict maximum, then $\varphi(\mathbf{r}) < \varphi(\mathbf{r}_0)$ for all $\mathbf{r}$ on $\partial B_R$ (for sufficiently small $R$), contradicting the equality. Therefore, $\varphi$ cannot have an interior maximum. The same argument applies to minima by considering $-\varphi$.

\textbf{Step 5: Establish uniqueness of solutions.}

Suppose $\varphi_1$ and $\varphi_2$ both satisfy $\nabla^{2}\varphi = 0$ with the same boundary conditions. Let $u = \varphi_1 - \varphi_2$. Then $\nabla^{2}u = 0$ with $u = 0$ on the boundary. Integrating $\nabla\cdot(u\nabla u)$ over the domain and using the divergence theorem:
\begin{equation}
\int_{\Omega} \nabla\cdot(u\nabla u)\,dV = \oint_{\partial\Omega} u\,\nabla u\cdot\hat{n}\,dA = 0.
\end{equation}
Since $\nabla\cdot(u\nabla u) = |\nabla u|^{2} + u\,\nabla^{2}u = |\nabla u|^{2}$, we obtain:
\begin{equation}
\int_{\Omega} |\nabla u|^{2}\,dV = 0 \quad\Longrightarrow\quad \nabla u = 0 \quad\Longrightarrow\quad u = \text{const} = 0.
\end{equation}
Hence $\varphi_1 = \varphi_2$ throughout the domain, proving uniqueness.

\textbf{Step 6: Write the general solution in spherical coordinates.}

Separating variables $\varphi(r,\theta,\phi) = R(r)\,\Theta(\theta)\,\Phi(\phi)$ in spherical coordinates yields Legendre's equation for $\Theta$ and Bessel-type solutions for $R$:
\begin{equation}
\varphi(r,\theta) = \sum_{\ell=0}^{\infty}\left(A_{\ell}\,r^{\ell} + B_{\ell}\,r^{-(\ell+1)}\right)P_{\ell}(\cos\theta),
\end{equation}
where $P_{\ell}$ are the Legendre polynomials. The coefficients $A_{\ell}$ and $B_{\ell}$ are determined by the boundary conditions. This is the multipole expansion of the potential.

\section*{Characteristics of the Equation}

\begin{itemize}
  \item \textbf{Maximum principle:} A harmonic function cannot have a local maximum or minimum in the interior of its domain; extrema occur on the boundary.
  \item \textbf{Mean-value property:} The value of $\varphi$ at any point equals the average of $\varphi$ over any sphere centered at that point.
  \item \textbf{Uniqueness:} For a given set of boundary conditions, the solution to Laplace's equation is unique.
  \item \textbf{Superposition:} Linear combinations of harmonic functions are harmonic.
  \item \textbf{Conformal invariance:} In 2D, $\nabla^{2}\varphi = 0$ is invariant under conformal (angle-preserving) mappings, which is the basis of conformal mapping techniques.
\end{itemize}
\section*{Solution Methods}

\begin{enumerate}
  \item \textbf{Separation of variables:} In Cartesian coordinates, assume $\varphi(x,y,z) = X(x)Y(y)Z(z)$; in spherical coordinates, $\varphi(r,\theta,\phi) = R(r)\Theta(\theta)\Phi(\phi)$, leading to Legendre polynomials and spherical harmonics $Y_{\ell m}(\theta,\phi)$; in cylindrical coordinates, Bessel functions appear.
  \item \textbf{Conformal mapping:} In 2D, any analytic function $f(z) = u + iv$ has harmonic real and imaginary parts. Mapping complicated geometries to simpler ones (e.g., via the Schwarz--Christoffel transformation) solves boundary value problems.
  \item \textbf{Green's function:} The solution is $\varphi(\mathbf{r}) = \oint G(\mathbf{r},\mathbf{r}')\,\sigma(\mathbf{r}')\,dA'$, where $G$ is the Green's function for the domain.
  \item \textbf{Numerical:} Finite difference, finite element, or boundary element methods solve $\nabla^{2}\varphi = 0$ on a discrete grid.
\end{enumerate}
\section*{Verifications}

Solutions to Laplace's equation are verified through experimental measurements of electric fields, temperature distributions, and fluid flow patterns. The capacitance of a spherical capacitor computed from $\varphi = Q/(4\pi\varepsilon_{0}r)$ agrees with measurement. Conformal mapping solutions for the flow around a cylinder match experimental streamlines in low-Reynolds-number flows.
\section*{Applications}

\begin{itemize}
  \item \textbf{Electrostatics:} The potential around conductors, in waveguides, and in capacitor geometries.
  \item \textbf{Gravitation:} The gravitational potential outside mass distributions; the geoid of the Earth.
  \item \textbf{Heat conduction:} Steady-state temperature distributions in solids.
  \item \textbf{Fluid mechanics:} Velocity potential for irrotational flow; flow around obstacles (potential flow theory).
  \item \textbf{Electrostatic shielding:} A conductor in electrostatic equilibrium is an equipotential; the field inside is zero---a consequence of the maximum principle.
\end{itemize}
\section*{What Richard Feynman Would Have Asked}

\subsection*{Plain-English Translation}
\textit{``What does Laplace's equation say, if I explain it without any math?''}

It says: in a region with no sources, the value of the potential at any point is exactly the average of its values on the surrounding boundary. If the boundary is hot, the interior must be hot; if the boundary is cold, the interior must be cold; and the temperature at any interior point is a weighted average of the boundary temperatures. This averaging property means the solution is always smooth---no bumps, no spikes, no surprises. The potential is as ``gentle'' as it can possibly be, subject to the boundary conditions.

\subsection*{Localized Mechanics}
\textit{``At the level of a small patch of space, what does $\nabla^{2}\varphi = 0$ physically mean?''}

The Laplacian $\nabla^{2}\varphi$ measures the difference between the value of $\varphi$ at a point and its average over a tiny surrounding sphere. If $\nabla^{2}\varphi > 0$, the point is a local ``valley'' (the average is higher); if $\nabla^{2}\varphi < 0$, it is a local ``peak'' (the average is lower). Setting $\nabla^{2}\varphi = 0$ means there are no local peaks or valleys: every point is exactly average. This is why harmonic functions are smooth and why they obey the maximum principle. Physically, in electrostatics, $\nabla^{2}\varphi = 0$ says that the electric flux into any small volume equals the flux out---charge conservation in the absence of sources. In heat conduction, it says the heat flowing into any small region equals the heat flowing out---thermal equilibrium.

\subsection*{Testing Extreme Limits}
\textit{``What happens to solutions of Laplace's equation in extreme geometries---very thin slits, sharp corners, or infinite domains?''}

Near a sharp corner of interior angle $\alpha$, the potential behaves as $r^{\pi/\alpha}$ (for a conducting corner). If $\alpha < \pi$ (a reentrant corner), the exponent $\pi/\alpha > 1$ and the electric field diverges as $r^{\pi/\alpha - 1} \to \infty$---this is the well-known ``lightning rod effect'' that explains why charge concentrates at sharp points. For an infinitely thin slit ($\alpha = 2\pi$), the field diverges as $r^{-1/2}$. In an infinite half-space with a boundary condition on the plane, the solution is obtained by the method of images: a point charge above a grounded plane has an image charge below, giving $\varphi = 0$ on the plane. In an infinite domain with no sources, the only bounded harmonic function is a constant (Liouville's theorem for harmonic functions). This is why the gravitational potential outside a spherically symmetric mass is uniquely $-GM/r$: it is the only bounded harmonic function that vanishes at infinity and has the correct monopole moment.

\subsection*{Dimensionality and Geometry}
\textit{``How does the dimensionality of space affect the solutions of Laplace's equation?''}

In 1D, Laplace's equation is $\varphi'' = 0$, with solutions $\varphi = ax + b$ (linear functions). In 2D, the fundamental solution is $\varphi = \ln r$ (logarithmic), and conformal mapping provides a powerful solution technique. In 3D, the fundamental solution is $\varphi = 1/r$ (Coulomb potential), and spherical harmonics provide the basis for separation of variables. In higher dimensions ($d \geq 3$), the fundamental solution is $\varphi = r^{2-d}$, which decays faster for larger $d$---gravity is weaker in higher dimensions. The change from $\ln r$ (2D) to $1/r$ (3D) reflects the different geometry of wavefronts: in 2D, wavefronts are circles ($\propto r$), while in 3D, they are spheres ($\propto r^{2}$). The density of states and the behavior of Green's functions are fundamentally different, leading to different physical consequences (e.g., bound states exist in 1D and 2D for arbitrarily weak potentials, but not in 3D---this is the threshold theorem for quantum mechanics).

\subsection*{Cross-Disciplinary Connection}
\textit{``Where does Laplace's equation or its solutions appear outside of classical physics?''}

In machine learning, the Laplacian operator appears in graph-based semi-supervised learning: the label function on a graph is ``harmonic'' (satisfies a discrete Laplace equation) on unlabeled nodes, with labeled nodes as boundary conditions. In image processing, Laplacian smoothing (diffusion of pixel values) is equivalent to solving the heat equation, whose steady state is Laplace's equation. In finance, the Black--Scholes equation can be transformed into Laplace's equation by a change of variables, and option pricing uses harmonic functions. In network theory, the effective resistance between two nodes in a resistor network satisfies a discrete Laplace equation (Kirchhoff's laws). In machine learning, the reproducing kernel Hilbert space associated with the Laplacian kernel $K(x,y) = \|x-y\|^{2-d}$ is used for interpolation. In all these fields, the key property exploited is the same: harmonic functions are determined by their boundary values.


% ============================================================
% BATCH 5 — Chapters 49–60
% Poynting Theorem through Stefan–Boltzmann Law
% ============================================================

% ------------------------------------------------------------
\section*{FAQ}

\begin{enumerate}
\item \textbf{Why does the maximum principle hold?}
If $\varphi$ had a local maximum at an interior point, then at least one second derivative would be negative there, but since all second derivatives sum to zero, at least one would have to be positive---contradicting the maximum. Physically, a potential cannot have an interior maximum because that would require a local concentration of ``nothing'' (no source), which is impossible.

\item \textbf{How does Laplace's equation differ from Poisson's equation?}
Poisson's equation $\nabla^{2}\varphi = -\rho/\varepsilon_{0}$ includes source terms. Laplace's equation is the special case $\rho = 0$. The general solution to Poisson's equation is the sum of a particular solution (from the sources) and a homogeneous solution (satisfying Laplace's equation, determined by boundary conditions).

\item \textbf{Is Laplace's equation linear?}
Yes. If $\varphi_{1}$ and $\varphi_{2}$ are solutions, so is $\alpha\varphi_{1} + \beta\varphi_{2}$. This linearity is why separation of variables and superposition work.
\end{enumerate}
\section*{Important Bibliography}

\begin{enumerate}
\item Jackson, J.D., \textit{Classical Electrodynamics}, 3rd ed. (Wiley, 1999), Chapters 1--3.
\item Haberman, R., \textit{Applied Partial Differential Equations with Fourier Series and Boundary Value Problems}, 5th ed. (Pearson, 2012), Chapter 10.
\item Arfken, G.B., Weber, H.J., and Harris, F.E., \textit{Mathematical Methods for Physicists}, 7th ed. (Academic Press, 2012), Chapter 17.
\item Zill, D.G., \textit{Advanced Engineering Mathematics}, 7th ed. (Jones \& Bartlett, 2021), Chapter 12.
\end{enumerate}

